\(\beta=1\), sample-specific-threshold member of the
taper family. Our mean-token-surprisal Countdown implementation is instead a
length-normalized analogue. Applying the same density-ratio coefficient to
Gaussian policies is an extension under a common dominating measure, not a
claim about the scope of the original TOPR analysis. Algebraically,
\begin{align}
e^{-[-\log\pi+\log\mu]_+}
&=\min\{\pi/\mu,1\},
\label{eq:topr_exact_identity}\\
e^{-[-L^{-1}\log\pi+L^{-1}\log\mu]_+}
&=\min\{\pi/\mu,1\}^{1/L}.
\label{eq:topr_length_normalized}
\end{align}
The Gaussian extension assumes a common dominating measure,
\(\mu(a_i)>0\), and finite support or integrable tails; using a direct
log-density ratio avoids conflating sample thresholds with unequal
normalization constants.

For a fixed reused sample, the far-branch conditional vector field contains
a \(\mu^{-1}\) coefficient, while the resulting absolute probability has a \(\mu\)
prefactor and the normalized ratio \(\pi/\mu\) decays as \(t^{-1}\). In the population
TOPR field, the outer sampling mass \(\mu\) cancels this denominator. We therefore
report weighted aggregate constants first and label any \(\mu_{\min}\) expression as
a finite-support worst-case relaxation. In particular, under the applicable
plateau branch,
\[
B_\omega^{\mathrm{TOPR}}
\le\sqrt{\lambda_{\max}(\Sigma^{-1})}
\sum_iq_i\sqrt{2\{\log(1/\mu_i)-C_\Sigma\}};
\]
otherwise the full branch-wise bound in
Appendix~\ref{app:gaussian_drpo_boundedness} is used.

\section{Experiments}
\label{sec:experiments}

The experiments proceed from occurrence to identification, control, and
transfer. We
first test whether the predicted far-field pattern appears in realistic
continuous-control and autoregressive policies. Controlled environments then
separate coefficient magnitude from policy geometry, trace the resulting
update into policy behavior, and compare negative-weighting rules under known
near/far structure. The final task-level sections retain the D4RL and Countdown
evaluation targets, including fields whose formal runs have not yet closed.

\subsection{Environments and Datasets}
\label{sec:environments_datasets}

We use two controlled environments and two external diagnostic domains. C-U1
is a contextual bandit with a Gaussian policy over continuous actions, while
D-U1 uses unordered categorical actions with shared representations. They
provide known ground truth for source isolation, policy dynamics, and
held-out-context generalization. D4RL Hopper and Countdown arithmetic
generation provide continuous offline-control and autoregressive sequence
diagnostics, respectively. Full configurations are in
Appendix~\ref{app:experimental_details}.

\paragraph{Measurement conventions.}
When the policy coordinate is controlled, source claims use policy-score
response. Neural diagnostics instead report the implemented actor-gradient
magnitude only to test whether the policy-level pattern appears in trainable
updates; such gradients also contain the network Jacobian.

\paragraph{Outcome taxonomy.}
We report task-performance collapse, policy-support or variance-boundary
events, and NaN/Inf numerical failure as distinct outcome classes. Their
protocol-specific definitions are in
Appendix~\ref{app:controlled_protocols}.

\subsection{External Diagnostic of Far-Field Repulsion}
\label{sec:external_far_field}

\begin{figure}[t]
    \centering
    \includegraphics[
        width=\columnwidth
    ]{figures/figure1_external_gradient_bottom_labels.pdf}

    \caption{
        \textbf{External diagnostic of far-field negative-update magnitude.}
        \textbf{(a)} Hopper: relative implemented actor-gradient diagnostic
        and matched absolute advantage as functions of relative standardized
        distance for transitions with negative learned-critic advantages.
        Values are normalized to the near-field reference within each seed.
        Lines report averages across 10 seeds, and shaded regions denote
        seed-level bootstrap 95\% confidence intervals.
        \textbf{(b)} Countdown: relative implemented actor-gradient diagnostic
        and fixed negative coefficient magnitude as functions of current-policy
        mean-token surprisal deciles among verifier-matched unsuccessful
        responses. Values are normalized to the lowest-surprisal decile, and
        shaded regions denote puzzle-level bootstrap 95\% confidence intervals.
    }
    \label{fig:external_far_field}
\end{figure}

Hopper provides a continuous-control setting with a state-conditioned Gaussian
policy,
where remoteness is measured by standardized distance from the current
policy. Countdown provides an autoregressive discrete setting, where
remoteness is measured by current-policy mean-token surprisal.

The diagnostic controls the sample-quality coefficient in each setting.
In Hopper, we compare negative-advantage transitions after matching their
absolute advantage magnitudes. In Countdown, we compare verifier-matched
unsuccessful responses while holding the negative coefficient magnitude
fixed. These matches test whether remoteness is reflected in implemented actor
updates independently of the sample-quality coefficient.

Figure~\ref{fig:external_far_field} shows signatures specific to each policy
family. In Hopper, the implemented actor-gradient diagnostic
increases with standardized distance while matched absolute advantage
remains approximately flat. In Countdown, the fixed-coefficient
diagnostic remains substantial at high mean-token surprisal. We interpret
the latter as persistent suppression in a categorical sequence policy,
not as an unbounded direct-logit score; implemented parameter gradients
also contain the network Jacobian. These diagnostics show that
remoteness-correlated update magnitude appears in realistic neural actors, but
they do not identify its source or causal effect.




\subsection{Controlled Identification of the Causal Mechanism}
\label{sec:controlled_identification}

The external diagnostics still entangle sample quality, remoteness, critic or
verifier coefficients, data coverage, and subsequent dynamics. Controlled
environments vary coefficient magnitude and remoteness independently and
intervene directly on near- and far-field negative updates. Full constructions
and protocols are provided in
Appendix~\ref{app:controlled_protocols}.

\subsubsection{Remoteness as an Independent Source of Update Magnitude}
\label{sec:controlled_source}

\begin{figure}[t]
    \centering

    \includegraphics[
        width=0.92\columnwidth,
        trim={0pt 2pt 0pt 2pt},
        clip
    ]{figures/fig_6_3_1_source_heatmap.pdf}

    \vspace{0.35em}

    \includegraphics[
        width=0.92\columnwidth,
        trim={0pt 2pt 0pt 2pt},
        clip
    ]{figures/fig_6_3_2_rescue_plot.pdf}

    \vspace{-0.4em}

\caption{
    \textbf{Controlled source isolation and causal transmission.}
    \textbf{Top:} source-isolation factorization across
    learner-relative remoteness bins, normalized to the near bin.
    COE denotes coefficient magnitude and SC denotes policy-score
    response; C rows report the C-U1 continuous protocol, and D rows
    show the corresponding direct-softmax categorical reference.
    \textbf{Bottom:} C-U1 causal-intervention summary. Ret. denotes
    reward retention; points show 20-seed means with bootstrap
    confidence intervals, and Task coll. reports task-performance
    collapse counts.
}
    \label{fig:controlled_6_3}
\end{figure}

We first isolate whether remoteness changes update magnitude independently of
sample quality. In C-U1, negative actions are
placed on the same reward contour, giving them identical fixed
advantages but different standardized distances from the current
Gaussian policy. The categorical direct-softmax reference applies the
same factorization by fixing the negative coefficient while varying
current-policy surprisal.

The top panel of Figure~\ref{fig:controlled_6_3} separates the
coefficient term from the policy-score term. The coefficient magnitude
remains flat across remoteness bins, whereas the score response increases
with remoteness. Thus, the larger update cannot be explained
by farther samples having larger advantages; it is generated by the
policy-score geometry itself. The manifestation differs across action
spaces: Gaussian policies exhibit amplitude amplification, while
categorical policies approach a bounded direct-logit score but continue
to suppress the selected action toward the simplex boundary. This
experiment identifies the source of the update difference, not its downstream
causal effect.

\subsubsection{Causal Transmission into Drift and Collapse}
\label{sec:controlled_transmission}

We next test whether the far-field update causally propagates into
instability. In the C-U1 intervention protocol, all branches share the
same data, fixed labels, initialization, optimizer, and training horizon,
and differ only in how near- and far-field negative updates are retained,
removed, capped, or globally rescaled. These matched-control invariants are
formalized in Appendix~\ref{app:controlled_shared}.

The bottom panel of Figure~\ref{fig:controlled_6_3} shows a clear rescue
pattern. Removing near-field negative updates leaves performance collapse
largely unchanged, whereas removing or capping far-field updates restores
reward retention and eliminates task-performance collapse. Global
rescaling also rescues the run, supporting the interpretation that
abnormal negative-update magnitude is a mediator. In contrast,
transferring the far-field budget to near-field negatives gives only
partial recovery, showing that the rescue is not obtained by simply
preserving the same total negative-update budget.

These interventions identify the far-field negative update, rather than
negative feedback in general, as a causal transmission path to policy
instability in the controlled Gaussian environment. Outcomes follow the
taxonomy in Section~\ref{sec:environments_datasets}.

\subsection{From Useful Repulsion to DRPO Control}
\label{sec:repulsion_control}

The causal rescue does not imply that all negative feedback should be removed.
C-U1 and D-U1 place a hidden optimum beyond the positive-only target and make
local repulsion point toward it, while additional remote negatives create
far-field pressure. This construction separates useful repulsion from harmful
dominance and permits controlled comparisons of taper shape.

\subsubsection{Stable Extrapolation from Controlled Negative Feedback}
\label{sec:controlled_negative_strength}

\begin{figure}[t]
    \centering
    \includegraphics[
        width=0.95\columnwidth
    ]{figures/fig_6_4_1_phase_transition.pdf}

    \vspace{-0.35em}
\caption{
    \textbf{Stable extrapolation from controlled negative feedback.}
    Increasing the effective negative strength $q/p$ first improves
    held-out-context reward beyond the positive-only ceiling, then causes
    over-extrapolation and task-performance collapse. Policy shift
    measures displacement from the positive-only target. The color strip
    reports task collapse, boundary events, and NaN/Inf failures
    separately.
}
    \label{fig:phase_transition_6_4_1}
\end{figure}

We sweep the effective negative strength relative to positive attraction,
denoted by $q/p$, under the matched-control invariants above; protocol and
metric details are in Appendix~\ref{app:negative_strength_sweep}. The
positive-only setting is $q/p=0$.

Figure~\ref{fig:phase_transition_6_4_1} shows a non-monotone transition.
For small and moderate $q/p$, held-out-context reward rises above the
positive-only ceiling, showing that controlled negative feedback provides
useful repulsion beyond the positive-support solution. As $q/p$ grows
larger, policy shift continues to increase but reward declines,
indicating over-extrapolation rather than additional useful
generalization. At the largest strengths, task-performance collapse and
boundary events appear; numerical failures remain a separate outcome class.

This strength sweep supplies the missing link between causal rescue and
tapered control: positive-only learning removes the harmful tail but also
removes the beneficial intermediate regime. Remoteness-aware control
should therefore preserve useful local repulsion while preventing the
far-field tail from dominating.


\begin{figure}[t]
\centering

\begin{minipage}[t]{0.68\columnwidth}
    \centering
    \vspace{0pt}
    \textbf{(a) Controlled taper comparison}\\[0.05em]
    \includegraphics[
        width=\linewidth
    ]{figures/fig_6_4_2_leftfig_bigtext_legend_protocol.pdf}
\end{minipage}
\hfill
\begin{minipage}[t]{0.28\columnwidth}
    \centering
    \vspace{0pt}
    \textbf{(b) Transfer}\\[0.20em]
    \scriptsize
    \setlength{\tabcolsep}{2.2pt}
    \renewcommand{\arraystretch}{0.92}
    \resizebox{\linewidth}{!}{
    \begin{tabular}{lcc}
    \toprule
    Method & Hop. & Cnt. \\
    \midrule
    Pos.        & \texttt{--} & \texttt{--} \\
    G-$\alpha$  & \texttt{--} & \texttt{--} \\
    Rec.-L      & \texttt{--} & \texttt{--} \\
    Rec.-Q      & \texttt{--} & \texttt{--} \\
    DRPO        & \textbf{\texttt{--}} & \textbf{\texttt{--}} \\
    \bottomrule
    \end{tabular}
    }

    \vspace{0.25em}
    {\scriptsize
    Hop.: Hopper return.\\
    Cnt.: Countdown success.
    }
\end{minipage}

\vspace{-0.35em}
\caption{
\textbf{Controlled taper comparison and external transfer.}
\textbf{(a)} Matched operating-point comparison in the controlled taper
protocol. For each method, the left bar is the near-field-retention matched
operating point and the right bar is the aggregate-budget matched operating
point. Bars decompose retained negative-update magnitude into near- and
far-field components, and dashed traces report held-out-context reward.
\textbf{(b)} Reserved external validation summary for Hopper and Countdown;
dashes mark outcomes awaiting the corresponding formal runs.
}
\label{fig:taper_control_and_transfer}
\end{figure}

\subsubsection{Controlled Comparison of Taper Shapes}
\label{sec:controlled_tapers}

We compare positive-only learning, global negative scaling,
reciprocal-linear tapering, reciprocal-quadratic tapering, and
exponential tapering while changing only the negative weighting rule. Two
matched comparisons are used. Near-field-retention matching preserves
comparable local negative-update magnitude and isolates differences in far-field
decay, whereas aggregate-budget matching tests whether a selective taper
outperforms a global reduction with the same total negative-update
budget.

As shown in Figure~\ref{fig:taper_control_and_transfer}, DRPO's exponential taper
provides the strongest stability--utility trade-off under both matched
comparisons. Positive-only learning avoids far-field instability but
remains at the positive-support ceiling. Global scaling suppresses useful
near- and harmful far-field feedback indiscriminately. Reciprocal
tapering introduces remoteness-aware suppression, but its tail remains
heavier